\par\medskip\noindent
2.5. \underline{Estructura de la categoría
$\operatorname{BIEXT}(P,Q;G)$}.

Las consideraciones formales generales de [4, n.º~1] proporcionan
varias consecuencias útiles de la dependencia aditiva de las
categorías $\operatorname{BIEXT}(P,Q;G)$ respecto de $G$.

2.5.1. En primer lugar, para $G$ fijo, la categoría
$\operatorname{BIEXT}(P,Q;G)$ está provista de una ley de composición
que asigna funtorialmente a dos biextensiones $E,E'$ de $(P,Q)$ por
$G$ una \underline{biextensión producto}, denotada $E.E'$. Esta
construcción no es sino el análogo del producto usual de dos
extensiones de un grupo conmutativo por otro. Es compatible con esa
construcción tanto si $E,E'$ se consideran extensiones de grupos de
$Q_P$ por $G_P$ como si se consideran extensiones de $P_Q$ por $G_Q$.
En particular, como torsor bajo $G_{P\times Q}$, $EE'$ es el producto

%%%%%%%%%%  scan idx 177  |  folio 166  %%%%%%%%%%

de los torsores bajo $G_{P\times Q}$ definidos por $E$ y $E'$.

La formación de $E.E'$ es además \underline{asociativa y conmutativa}
salvo isomorfismos canónicos, y esta ley de composición sobre
$\operatorname{BIEXT}(P,Q;G)$ posee un objeto neutro bilateral,
llamado la \underline{biextensión trivial de $(P,Q)$ por $G$}. Puede
construirse, como en \emph{loc.\ cit.}, a partir de la única
biextensión $E_0$ de $(P,Q)$ por el grupo nulo $0$, salvo isomorfismo,
mediante el homomorfismo $0\longrightarrow G$. De esta descripción se
obtiene inmediatamente otra descripción de la biextensión trivial como
el torsor trivial
\[
 E=P\times Q\times G.
\]
Sus dos estructuras de grupo sobre $P$ y sobre $Q$ se obtienen,
respectivamente, de las estructuras de grupo producto sobre
$Q\times G$ y $P\times G$ por los cambios de base
$P\longrightarrow e$ y $Q\longrightarrow e$.

Finalmente, para toda biextensión $E$ de $(P,Q)$ por $G$, puede
definirse la \underline{biextensión inversa}, denotada $E^{-1}$, junto
con un isomorfismo de $EE^{-1}$ con la biextensión trivial. Esto
determina $E^{-1}$ salvo isomorfismo único y hace funtorial su
dependencia respecto de $E$.

En resumen, la categoría $\operatorname{BIEXT}(P,Q;G)$ se comporta
como un objeto grupo de la categoría $\mathbf{Cat}$ de categorías,
salvo que las relaciones de asociatividad, conmutatividad, unidad e
inversión solo se cumplen salvo isomorfismos especificados. Puede
llamarse un \underline{pseudogrupo} en la 2-categoría $\mathbf{Cat}$.

%%%%%%%%%%  scan idx 178  |  folio 167  %%%%%%%%%%

2.5.2. Para $P,Q,G$ fijos, las propiedades precedentes implican los
hechos siguientes. Para toda biextensión $E$ de $(P,Q)$ por $G$, el
monoide de endomorfismos de $E$ es canónicamente isomorfo al monoide de
endomorfismos de la biextensión trivial $\underline 1$: a un
endomorfismo $u$ de $\underline 1$ se le asocia
$u.\mathrm{id}_E$ sobre $\underline 1.E\simeq E$. El monoide de
endomorfismos de $\underline 1$ es, de hecho, un
\underline{grupo conmutativo}, denotado
\begin{equation}\tag{2.5.2.1}
 \operatorname{Biext}^{0}(P,Q;G).
\end{equation}
Su ley de composición viene dada también por
$(u,v)\mapsto u.v\in
\operatorname{End}(\underline 1.\underline 1)
\simeq\operatorname{End}(\underline 1)$.

Asimismo, la ley de composición $(E,E')\mapsto EE'$ dota al conjunto
\begin{equation}\tag{2.5.2.2}
 \operatorname{Biext}^{1}(P,Q;G)
\end{equation}
de clases de isomorfismo de biextensiones de $(P,Q)$ por $G$ de una
estructura de \underline{grupo conmutativo}; el inverso de la clase de
$E$ es la clase de $E^{-1}$.

2.5.3. Dado un homomorfismo de grupos $G\longrightarrow G'$, el funtor
correspondiente
\begin{equation}\tag{2.5.3.1}
 \operatorname{BIEXT}(P,Q;G)
 \longrightarrow
 \operatorname{BIEXT}(P,Q;G')
\end{equation}
es compatible, salvo isomorfismos canónicos, con las operaciones
$E.E'$ y $E^{-1}$ de 2.5.1. En consecuencia, las dos aplicaciones
evidentes
\begin{equation}\tag{2.5.3.2}
 \operatorname{Biext}^{i}(P,Q;G)
 \longrightarrow
 \operatorname{Biext}^{i}(P,Q;G')
 \qquad(i=0,1)
\end{equation}
son \underline{homomorfismos de grupos}. Además, si $E$ es una
biextensión de $(P,Q)$ por $G$ y $E'$ la biextensión correspondiente
por $G'$, el homomorfismo natural
\begin{equation}\tag{2.5.3.3}
 \operatorname{Hom}(E,E)=\operatorname{Aut}(E)
 \longrightarrow
 \operatorname{Hom}(E',E')=\operatorname{Aut}(E')
\end{equation}
es precisamente (2.5.3.2) para $i=0$.

%%%%%%%%%%  scan idx 179  |  folio 168  %%%%%%%%%%

2.5.4. El grupo $\operatorname{Biext}^{0}(P,Q;G)$ admite una
descripción explícita sencilla. Si $E$ es una biextensión de $(P,Q)$
por $G$, un automorfismo $u$ de $E$ es ante todo un automorfismo del
torsor subyacente bajo $G_{P\times Q}$ y, por tanto, viene dado por un
morfismo
\[
 u_0:P\times Q\longrightarrow G
\]
mediante la fórmula
\[
 u(x)=u_0(j(x))x.
\]
La compatibilidad de $u$ con la primera, respectivamente con la
segunda, estructura de extensión sobre $E$ equivale a la aditividad de
$u_0$ respecto del segundo, respectivamente del primer, argumento.
Así, \underline{$u$ es un automorfismo de la estructura de biextensión
de $E$ si y solo si $u_0$ es bilineal}. Se obtiene una
\underline{biyección canónica}
\begin{equation}\tag{2.5.4.1}
 \operatorname{Biext}^{0}(P,Q;G)
 \simeq
 \operatorname{Hom}(P\otimes Q,G).
\end{equation}
Se comprueba fácilmente que esta biyección es independiente de $E$,
que respeta las estructuras de grupo y que es funtorial en $G$ para la
funtorialidad del primer miembro descrita en (2.5.3.2).

Más adelante determinaremos $\operatorname{Biext}^{1}(P,Q;G)$ y
veremos que es canónica y funtorialmente en $G$ isomorfo a
\[
 \operatorname{Ext}^{1}
 \bigl(P\overset{\mathbb L}{\otimes}Q,G\bigr).
\]
